\chapter{联络、曲率与度量几何}
\label{chap:math-connections}

上一章只定义了每个点的切空间。不同点的切向量属于不同向量空间，表达式
\(Y(q)-Y(p)\) 没有类型意义。联络的职责不是“把普通导数写弯”，而是补充一个
跨点比较规则。

\begin{definition}[向量丛与截面]
向量丛 \(E\xrightarrow{\pi}M\) 给每个 \(p\in M\) 配一个同维向量空间
\(E_p=\pi^{-1}(p)\)，并要求在每个足够小的开集 \(U\) 上能光滑地识别
\(\pi^{-1}(U)\) 与 \(U\times\mathbb R^r\)。
一个截面 \(s\) 给每个 \(p\) 选取 \(s(p)\in E_p\)，且随点光滑变化；
所有截面的集合记为 \(\Gamma(E)\)。切丛 \(TM\) 是 \(E_p=T_pM\)
的特例。
\end{definition}

局部识别允许在一小片区域中用 \(r\) 个函数记录截面，但不同局部识别的
基一般随位置变化。若直接对这些函数求导，结果会依赖选用的基。
联络补偿这种基的变化。

\begin{definition}[向量丛上的联络]
向量丛 \(E\to M\) 上的联络是映射
\[
\nabla:\Gamma(TM)\times\Gamma(E)\to\Gamma(E)
\]
满足对第一变量 \(C^\infty(M)\)-线性
（这里 \(C^\infty(M)\) 是流形上的光滑实值函数），
对第二变量 \(\mathbb R\)-线性，并满足
\[
\nabla_X(f s)=X(f)s+f\nabla_Xs.
\]
\end{definition}

\section{联络系数从定义推出}
\label{sec:math-connection-coefficients}

对切丛 \(E=TM\)，在坐标基 \(\partial_i\) 上定义函数
\(\Gamma^k{}_{ij}\)：
\[
\nabla_{\partial_i}\partial_j=\Gamma^k{}_{ij}\partial_k.
\]
若 \(X=X^i\partial_i\)、\(Y=Y^j\partial_j\)，联络对第一个变量的
函数线性性及对第二个变量的乘积法则给
\begin{equation}
(\nabla_XY)^k
=X^i\bigl(\partial_iY^k+\Gamma^k{}_{ij}Y^j\bigr).
\label{eq:math-covariant-derivative-coordinates}
\end{equation}
第一个加数对坐标分量求导，第二个加数补偿坐标基随位置的变化。

为什么 \(\Gamma^k{}_{ij}\) 不是张量？令新坐标为
\(y^a=y^a(x)\)。前章已证明
\(\partial/\partial y^a=(\partial x^i/\partial y^a)\partial_i\)。
将此式代入联络定义，先对第二个基的系数用乘积法则，再把输出基
换成 \(\partial/\partial y^c\)，得到
\begin{equation}
\Gamma'^c{}_{ab}
=\frac{\partial y^c}{\partial x^k}
\left(
\frac{\partial x^i}{\partial y^a}
\frac{\partial x^j}{\partial y^b}\Gamma^k{}_{ij}
+\frac{\partial^2x^k}{\partial y^a\partial y^b}
\right).
\label{eq:math-connection-transform}
\end{equation}
括号最后一项含坐标变换的二阶导数，张量的变换律没有这种项。
把式~\eqref{eq:math-connection-transform} 代回
式~\eqref{eq:math-covariant-derivative-coordinates}，它恰消掉
\(\partial_iY^k\) 变换时出现的多余项，使 \(\nabla_XY\) 真正按向量变换。

平面是检验该式的简单反例。在 Cartesian 坐标中取通常平直联络，
\(\Gamma^k{}_{ij}=0\)。改用极坐标 \((r,\theta)\)，直接对
\(\mathbf e_r=(\cos\theta,\sin\theta)\)、
\(\partial_\theta=r\mathbf e_\theta\) 求导，得到
\[
\nabla_{\partial_\theta}\partial_\theta=-r\partial_r,\qquad
\nabla_{\partial_r}\partial_\theta
=\nabla_{\partial_\theta}\partial_r=r^{-1}\partial_\theta.
\]
故 \(\Gamma^r{}_{\theta\theta}=-r\)、
\(\Gamma^\theta{}_{r\theta}=\Gamma^\theta{}_{\theta r}=1/r\)，
即使平面本身没有弯曲，联络系数也可以非零。

\section{平行移动是可解的初值问题}

给定曲线 \(\gamma(t)\) 与沿曲线的向量 \(V(t)=V^k(t)\partial_k\)，
把它代入式~\eqref{eq:math-covariant-derivative-coordinates}，得到
\[
\nabla_{\dot\gamma}V=0
\quad\Longleftrightarrow\quad
\dot V^k+\Gamma^k{}_{ij}(\gamma(t))
\dot\gamma^i V^j=0.
\]
右边是一组线性常微分方程。若系数沿曲线连续，给定
\(V(t_0)\in T_{\gamma(t_0)}M\) 后，在曲线存在的区间中局部有唯一解。
这样联络把起点的向量搬到邻近点，且无需先把两个切空间直接相减。

若联络还满足度量相容条件
\(X[g(Y,Z)]=g(\nabla_XY,Z)+g(Y,\nabla_XZ)\)，
则两个平行向量 \(V,W\) 沿同一曲线满足
\[
\frac{\dd}{\dd t}g(V,W)
=g(\nabla_{\dot\gamma}V,W)+g(V,\nabla_{\dot\gamma}W)=0.
\]
也就是说，长度和夹角在平行移动中保持；下一节将问什么条件使这样的
联络由度量唯一决定。

把前面的平面极坐标例子算到底，能避免一种常见误解。
沿半径为 \(r_0\) 的圆周取 \(\gamma(t)=(r_0,t)\)，
\(0\le t\le2\pi\)。平行方程成为
\[
\dot V^r-r_0V^\theta=0,
\qquad
\dot V^\theta+r_0^{-1}V^r=0.
\]
第二个坐标基向量 \(\partial_\theta\) 的长度是 \(r_0\)，
所以令 \(W^\theta=r_0V^\theta\) 才得到正交基中的物理分量。
此时 \(\dot V^r=W^\theta\)、\(\dot W^\theta=-V^r\)。
若起点分量为 \((A,B)\)，解为
\[
V^r(t)=A\cos t+B\sin t,
\qquad
W^\theta(t)=-A\sin t+B\cos t.
\]
走满一圈后分量回到 \((A,B)\)。沿途联络系数非零，
但闭路平行移动的净结果仍是恒等；这与平面曲率为零一致。

\begin{exercise}
在同一平面取不是闭合圆周的弧线 \(0\le t\le\pi/2\)，
将上式的正交基分量与 Cartesian 分量比较。
解释为什么正交基分量旋转了 \(\pi/2\)，平行移动的向量
在固定 Cartesian 坐标中却没有旋转。
\end{exercise}


\begin{definition}[伪 Riemann 度量]
度量是每点切空间上的非退化对称双线性型 \(g_p\)，并随 \(p\) 光滑变化。
若正定，称 Riemann 度量；若签名不全正，称伪 Riemann 度量。
\end{definition}

度量并非某张坐标图中的固定矩阵。若 \(y=y(x)\)，
把前章的基变换代入同一双线性型，得到
\[
g'_{ab}=g\!\left(\frac{\partial}{\partial y^a},
\frac{\partial}{\partial y^b}\right)
=\frac{\partial x^i}{\partial y^a}
\frac{\partial x^j}{\partial y^b}g_{ij}.
\]
例如 Euclid 平面的 \(g=dx^2+dy^2\)，在
\(x=r\cos\theta,y=r\sin\theta\) 下满足
\(dx=\cos\theta\,dr-r\sin\theta\,d\theta\)、
\(dy=\sin\theta\,dr+r\cos\theta\,d\theta\)。
两式平方相加时交叉项抵消，得
\(g=dr^2+r^2d\theta^2\)。
极坐标矩阵在 \(r=0\) 看似退化，是该坐标图在那里失效，
并非原平面度量突然失去非退化性。

\section{度量怎样唯一选出联络}
\label{sec:math-levi-civita-proof}

度量 \(g\) 量出同一点两向量的内积；联络告诉我们怎样跨点搬运向量。
仅要求平行移动保持内积还不能选出唯一联络，因为沿闭路搬运时仍可附加
局部旋转。再要求没有挠率，就把这份自由度固定了。

\begin{definition}[挠率与度量相容]
切丛联络的挠率是
\[
T(X,Y)=\nabla_XY-\nabla_YX-[X,Y].
\]
若 \(T=0\)，称联络无挠。若对任意向量场 \(X,Y,Z\) 有
\[
X[g(Y,Z)]
=g(\nabla_XY,Z)+g(Y,\nabla_XZ),
\]
称联络与度量 \(g\) 相容。
\end{definition}

在坐标基中，\([\partial_i,\partial_j]=0\)，故无挠等价于
\(\Gamma^k{}_{ij}=\Gamma^k{}_{ji}\)。度量相容则等价于
\(\partial_i g_{jk}
=\Gamma^\ell{}_{ij}g_{\ell k}
+\Gamma^\ell{}_{ik}g_{j\ell}\)。
这些是两组可直接求解联络系数的线性方程。

\begin{theorem}[Levi--Civita 联络]
对任意光滑、非退化的对称度量 \(g\)，存在唯一既与 \(g\) 相容又无挠的
切丛联络。它满足 Koszul 公式
\begin{align*}
2g(\nabla_XY,Z)
={}&Xg(Y,Z)+Yg(Z,X)-Zg(X,Y)\\
&+g([X,Y],Z)-g([Y,Z],X)-g([X,Z],Y),
\end{align*}
在坐标基中等价于
\begin{equation}
\Gamma^k{}_{ij}
=\frac12g^{k\ell}
(\partial_i g_{j\ell}+\partial_j g_{i\ell}
-\partial_\ell g_{ij}).
\label{eq:math-levi-civita-coefficients}
\end{equation}
\end{theorem}
\begin{proof}
先证唯一性。对 \(g(Y,Z),g(Z,X),g(X,Y)\) 分别沿
\(X,Y,Z\) 求导，用度量相容展开，再把前两式相加、第三式相减。
无挠性允许把 \(\nabla_YX\) 换成 \(\nabla_XY-[X,Y]\)，
把 \(\nabla_XZ\) 换成 \(\nabla_ZX+[X,Z]\)，
把 \(\nabla_YZ\) 换成 \(\nabla_ZY+[Y,Z]\)。
含 \(\nabla_ZX,\nabla_ZY\) 的项相消，留下
\(2g(\nabla_XY,Z)\) 和上述三个 Lie 括号项。
由于 \(g\) 非退化，它与所有 \(Z\) 的配对唯一决定 \(\nabla_XY\)。

再证存在性。把 Koszul 公式右边视为对 \(Z\) 的线性泛函，
由 \(g\) 非退化定义唯一向量 \(\nabla_XY\)。
将 \(Y\) 换成 \(fY\)，右边对 \(f\) 求导的项化为
\(2X(f)g(Y,Z)\)，于是
\(\nabla_X(fY)=X(f)Y+f\nabla_XY\)；
将 \(X\) 换成 \(fX\)，所有 \(Y(f),Z(f)\) 项两两相消，得到
\(\nabla_{fX}Y=f\nabla_XY\)。
交换 \(X,Y\) 后相减，右边给
\(2g([X,Y],Z)\)，故 \(T=0\)。
对 \(Y,Z\) 交换并相加，Lie 括号项抵消，得到度量相容式。
这验证构造的确是所需联络。
最后在 \(X=\partial_i,Y=\partial_j,Z=\partial_\ell\)
时代入 \([\partial_i,\partial_j]=0\)，再乘 \(g^{k\ell}\)，
得到坐标公式。
\end{proof}

定理中“非退化”不可删：若 \(g(v,\cdot)=0\) 可对非零 \(v\) 成立，
Koszul 公式就不能由与所有 \(Z\) 的配对唯一找回 \(\nabla_XY\)。
而度量签名不影响证明，所以 Lorentz 度量也有 Levi--Civita 联络。

\section{测地线的两种定义}
\label{sec:math-geodesic-variation}

平行移动已由线性常微分方程定义。一条曲线若自己的速度沿自身平行，
就称为仿射参数化的测地线：
\begin{equation}
\nabla_{\dot\gamma}\dot\gamma=0
\quad\Longleftrightarrow\quad
\ddot x^k+\Gamma^k{}_{ij}\dot x^i\dot x^j=0.
\label{eq:math-geodesic-equation}
\end{equation}
在正定度量中，这与能量泛函
\(E[\gamma]=\frac12\int_{t_0}^{t_1}
g_{ij}(x)\dot x^i\dot x^j\,\dd t\) 的固定端点驻值相同。
为看清联系，直接对 Lagrangian
\(L=\frac12g_{ij}\dot x^i\dot x^j\) 使用 Euler--Lagrange 方程：
\[
\frac{\dd}{\dd t}(g_{kj}\dot x^j)
-\frac12\partial_k g_{ij}\dot x^i\dot x^j=0.
\]
展开时间导数并利用 \(\dot x^i\dot x^j\) 的对称性，再乘 \(g^{\ell k}\)，
恰得到式~\eqref{eq:math-geodesic-equation}。
固定端点使分部积分的边界项消失；“驻值”并不保证整条测地线在全局最短。
Lorentz 度量的零测地线长度为零，仍可由仿射方程定义，
因此也不应称作“最短路”。

\begin{exercise}
把 \(g=dr^2+r^2d\theta^2\) 代入
式~\eqref{eq:math-levi-civita-coefficients}，
重新得到上一节的平面极坐标联络系数。
\end{exercise}


联络决定平行移动以后，还可以问：沿两个方向组成的小闭路搬运一个向量，
最后是否回到原向量？两个协变导数不对易的部分给出曲率。

\begin{proposition}[曲率的张量性]
\[
R(X,Y)Z=\nabla_X\nabla_YZ-\nabla_Y\nabla_XZ-\nabla_{[X,Y]}Z
\]
对 \(X,Y,Z\) 的每个变量都是 \(C^\infty(M)\)-线性的。
\end{proposition}
\begin{proof}
对第一变量，把 \(X\) 换成 \(fX\)。第二项
\(-\nabla_Y\nabla_{fX}Z\) 多出 \(-Y(f)\nabla_XZ\)，而
\([fX,Y]=f[X,Y]-Y(f)X\) 使第三项多出
\(+Y(f)\nabla_XZ\)，故 \(R(fX,Y)Z=fR(X,Y)Z\)。
由定义 \(R(X,Y)=-R(Y,X)\)，第二变量也具有线性性。
最后将 \(Z\) 换成 \(fZ\)，前两项的交叉导数互相抵消，
剩余的新项为 \((XY-YX)(f)Z\)；第三项恰减去
\([X,Y](f)Z\)。由于 Lie 括号定义为 \(XY-YX\)，余项为零。
\end{proof}

\section{曲率分量从两次协变求导得到}
\label{sec:math-curvature-components}

在坐标基中，\([\partial_i,\partial_j]=0\)。令
\(R(\partial_i,\partial_j)\partial_k
=R^\ell{}_{kij}\partial_\ell\)。
对 \(\nabla_{\partial_j}\partial_k
=\Gamma^\ell{}_{jk}\partial_\ell\) 再沿 \(\partial_i\) 求导，
分别对系数和基使用联络的乘积法则，便得
\begin{equation}
R^\ell{}_{kij}
=\partial_i\Gamma^\ell{}_{jk}
-\partial_j\Gamma^\ell{}_{ik}
+\Gamma^m{}_{jk}\Gamma^\ell{}_{im}
-\Gamma^m{}_{ik}\Gamma^\ell{}_{jm}.
\label{eq:math-curvature-components}
\end{equation}
每项都有来处：前两项是联络系数随位置的变化，后两项是
第一次求导后产生的基向量还要再次求导。

平面极坐标能检验“联络系数非零不代表曲率非零”。
上一节得到 \(\Gamma^r{}_{\theta\theta}=-r\)、
\(\Gamma^\theta{}_{r\theta}
=\Gamma^\theta{}_{\theta r}=1/r\)。
由\cref{eq:math-curvature-components}，
\[
R^r{}_{\theta r\theta}
=\partial_r(-r)
-\Gamma^\theta{}_{r\theta}\Gamma^r{}_{\theta\theta}
=-1-\frac1r(-r)=0.
\]
其余独立分量也为零；平面极坐标中的非零 \(\Gamma\)
仅反映基向量随位置改变。

\section{换用活动标架时怎样记录同一个曲率}

在非坐标标架 \(e_a\) 中，定义联络一形式
\(\omega^a{}_b(X)\) 使
\(\nabla_Xe_b=\omega^a{}_b(X)e_a\)。
把曲率定义作用于 \(e_b\)，使用一形式外微分的公式
\[
\dd\omega(X,Y)
=X[\omega(Y)]-Y[\omega(X)]-\omega([X,Y]),
\]
可整理为
\[
R(X,Y)e_b=\Omega^a{}_b(X,Y)e_a,\qquad
\Omega^a{}_b=\dd\omega^a{}_b+\omega^a{}_c\wedge\omega^c{}_b.
\]
第一项来自对联络系数求导，第二项来自连续两次基变换。
它与\cref{eq:math-curvature-components} 是同一张量的两种记法。

若 \(\theta^a\) 与 \(e_a\) 对偶，挠率形式
\(T^a(X,Y)=\theta^a(T(X,Y))\) 同样满足
\[
T^a=\dd\theta^a+\omega^a{}_b\wedge\theta^b.
\]
证明只需将右边作用于 \(e_b,e_c\)：\(\dd\theta^a\)
给出 \(-\theta^a([e_b,e_c])\)，另两项分别给出
\(\theta^a(\nabla_{e_b}e_c)\) 和
\(-\theta^a(\nabla_{e_c}e_b)\)，恰为挠率定义。
这说明标架不对易与几何挠率不是一回事：
Levi--Civita 联络虽无挠，非坐标标架仍可有 \(\dd\theta^a\ne0\)。

\begin{exercise}
把单位球面的度量整体乘 \(a^2\)，其中 \(a>0\)。
从新的正交余标架求 \(\omega^1{}_2\) 与 \(\Omega^1{}_2\)，
证明半径为 \(a\) 的球面 Gauss 曲率为 \(1/a^2\)。
\end{exercise}


\section{球面的贯穿计算}

单位球面坐标 \((\vartheta,\varphi)\) 给
\[
g_{\vartheta\vartheta}=1,\qquad
g_{\varphi\varphi}=\sin^2\vartheta .
\]
逆矩阵只有 \(g^{\vartheta\vartheta}=1\) 和
\(g^{\varphi\varphi}=\sin^{-2}\vartheta\) 两个非零分量。
代入 Levi--Civita 公式，所有非零系数为
\[
\Gamma^\vartheta_{\varphi\varphi}
=-\sin\vartheta\cos\vartheta,\qquad
\Gamma^\varphi_{\vartheta\varphi}
=\Gamma^\varphi_{\varphi\vartheta}=\cot\vartheta.
\]
按前述 \(R(X,Y)Z\) 的号差约定，选择
\(X=\partial_\vartheta,\ Y=\partial_\varphi,\ Z=\partial_\varphi\)。
在 \(\partial_\vartheta\) 方向的曲率分量中，真正留下的是
\[
\partial_\vartheta\Gamma^\vartheta_{\varphi\varphi}
-\Gamma^\vartheta_{\varphi\varphi}
 \Gamma^\varphi_{\vartheta\varphi}
=-(\cos^2\vartheta-\sin^2\vartheta)+\cos^2\vartheta
=\sin^2\vartheta .
\]
因而
\[
R^\vartheta{}_{\varphi\vartheta\varphi}=\sin^2\vartheta.
\]
二维曲面的 Gauss 曲率定义为
\(K=g(R(X,Y)Y,X)/
[g(X,X)g(Y,Y)-g(X,Y)^2]\)，其中 \(X,Y\) 线性独立。
分子和分母在换基时都乘同一个行列式平方，故商不依赖所选的
\(X,Y\)。在坐标基上，
\(R_{\vartheta\varphi\vartheta\varphi}
=g_{\vartheta\vartheta}R^\vartheta{}_{\varphi\vartheta\varphi}
=\sin^2\vartheta\)。
除以 \(\det(g_{ij})=\sin^2\vartheta\)，得到 Gauss 曲率 \(K=1\)。
分量随坐标而变，标量曲率不会；极点处的 \(\sin\vartheta=0\)
只说明这张经纬度图失效，不说明球面曲率发散。

同一球面还能把曲率解释成闭路平行移动的净转角。
在避开极点的区域取正交余标架
\(\theta^1=d\vartheta\)、\(\theta^2=\sin\vartheta\,d\varphi\)。
度量相容性使 \(\omega^2{}_1=-\omega^1{}_2\)；
无挠方程 \(d\theta^a+\omega^a{}_b\wedge\theta^b=0\)
与 \(d\theta^2=\cos\vartheta\,d\vartheta\wedge d\varphi\)
共同给出
\[
\omega^1{}_2=-\cos\vartheta\,d\varphi,
\qquad
\Omega^1{}_2=d\omega^1{}_2
=\sin\vartheta\,d\vartheta\wedge d\varphi
=\theta^1\wedge\theta^2.
\]
固定纬度 \(\vartheta=\vartheta_0\)，沿 \(\varphi:0\to2\pi\)
平行搬运 \(V=V^ae_a\)。分量方程是
\(dV^1/d\varphi=\cos\vartheta_0 V^2\)、
\(dV^2/d\varphi=-\cos\vartheta_0 V^1\)，
故回到起点时相对同一正交标架转过
\(-2\pi\cos\vartheta_0\)。转角只在模 \(2\pi\) 意义下确定，
它也可写成 \(2\pi(1-\cos\vartheta_0)\)，
正好是北极到该纬线所围球冠的面积。
与前面的平面圆周不同，这条闭路保留了非零净转角。

\section{曲率怎样改变相邻的自由轨迹}

仅把曲率当作四指标数组，还不能看出它对运动的影响。
设 \(\gamma_s(t)\) 是一族以仿射参数表示的测地线，
定义速度场 \(u=\partial_t\gamma_s\) 和连接相邻轨迹的
变分场 \(J=\partial_s\gamma_s\)。参数 \(s,t\) 可交换，
所以 \([J,u]=0\)。联络无挠给
\(\nabla_uJ=\nabla_Ju\)，每条曲线为测地线给
\(\nabla_uu=0\)。把这两式代入曲率定义，得到
\[
\nabla_u\nabla_uJ
=\nabla_u\nabla_Ju
=R(u,J)u+\nabla_J\nabla_uu
=-R(J,u)u.
\]
因此相邻测地线的间距满足 Jacobi 方程
\begin{equation}
\nabla_u^2J+R(J,u)u=0.
\label{eq:math-jacobi-equation}
\end{equation}
这不是在同一条轨迹上人为加入“曲率力”；
它比较的是初始条件略有差别的两条自由轨迹。
在单位球面取单位速的大圆，且 \(J\perp u\)，
前面算出的 \(K=1\) 意味着 \(R(J,u)u=J\)。
若沿测地线选平行单位法向量 \(n\)，写 \(J=j(t)n\)，
就得到 \(j''+j=0\)。从同一北极以略不同方向出发的
大圆在南极再次相遇，正是这个振荡解的几何含义。

\section{等距流、Killing 场与守恒量}

有了联络和度量，便能严格表达“空间沿某方向移动后看起来不变”。
设向量场 \(K\) 的局部流为 \(\Phi_t\)：从点 \(p\) 出发，
曲线 \(t\mapsto\Phi_t(p)\) 满足
\(\frac{\dd}{\dd t}\Phi_t(p)=K_{\Phi_t(p)}\)、\(\Phi_0(p)=p\)。
如果流把度量保持不变，即
\(g_{\Phi_t(p)}(\dd\Phi_t X,\dd\Phi_t Y)=g_p(X,Y)\)，
就称 \(\Phi_t\) 为局部等距流。

\begin{definition}[Lie 导数与 Killing 场]
度量沿 \(K\) 的 Lie 导数定义为
\[
(\mathcal L_Kg)_p(X,Y)
=\left.\frac{\dd}{\dd t}\right|_{t=0}
(\Phi_t^*g)_p(X,Y).
\]
若 \(\mathcal L_Kg=0\)，称 \(K\) 为 Killing 场。
\end{definition}

这里 \(\Phi_t^*g\) 表示先用微分 \(\dd\Phi_t\)
把 \(p\) 处两个向量推到 \(\Phi_t(p)\)，再在那里计算 \(g\)。
所以定义不是“直接微分不同点的矩阵”，而是先把不同点的双线性型
拉回同一切空间再比较。

\begin{proposition}[Killing 方程]
对 Levi--Civita 联络，\(K\) 是 Killing 场当且仅当
\[
g(\nabla_XK,Y)+g(X,\nabla_YK)=0
\quad\text{对所有 }X,Y.
\]
在坐标中这等价于
\(\nabla_\mu K_\nu+\nabla_\nu K_\mu=0\)。
\end{proposition}
\begin{proof}
从流的定义对 \(t\) 求导，或在一张坐标图中直接对
\(g_{kl}(\Phi_t(x))\partial_i\Phi_t^k\partial_j\Phi_t^l\)
求导，得到
\[
(\mathcal L_Kg)(X,Y)
=K[g(X,Y)]-g([K,X],Y)-g(X,[K,Y]).
\]
度量相容性给
\(K[g(X,Y)]=g(\nabla_KX,Y)+g(X,\nabla_KY)\)。
无挠性给 \([K,X]=\nabla_KX-\nabla_XK\)，对 \(Y\) 亦然。
代入并约去 \(\nabla_KX,\nabla_KY\) 项，便得到命题中的式子。
\end{proof}

若 \(\gamma(t)\) 是仿射参数测地线，令 \(u=\dot\gamma\)，
则 \(\nabla_uu=0\)。因此
\[
\frac{\dd}{\dd t}g(K,u)
=g(\nabla_uK,u)+g(K,\nabla_uu)=0.
\]
第一项是 Killing 方程取 \(X=Y=u\) 的一半，第二项由
测地线方程消失。这给出守恒量 \(g(K,u)\)。
例如 Euclid 平面中 \(K=\partial_x\) 产生平移，
\(g(K,u)=\dot x\) 守恒；\(K=-y\partial_x+x\partial_y\)
产生旋转，\(g(K,u)=x\dot y-y\dot x\) 守恒。
由此可见“对称性给守恒量”不是额外的口号，而是度量相容性
与 Killing 方程的直接推论。

\section{Hodge 对偶从内积与体积形式定义}

上一章已把 \(k\)-形式定义为交替的 \(k\) 重协变张量。
现在先取定向的 \(n\) 维 Riemann 流形，即度量正定；
伪 Riemann 情形的符号稍后单独说明。
在每点选正向正交余标架 \(e^1,\ldots,e^n\)。
体积形式是
\(\mathrm{vol}_g=e^1\wedge\cdots\wedge e^n\)。
若 \(I=(i_1<\cdots<i_k)\)，记
\(e^I=e^{i_1}\wedge\cdots\wedge e^{i_k}\)，
规定这些 \(e^I\) 为 \(\Lambda^kT_p^*M\) 的正交归一基；
这将向量上的度量唯一延拓为 \(k\)-形式上的内积。

\begin{definition}[Hodge 星号]
Hodge 算子 \(*:\Lambda^kT_p^*M\to\Lambda^{n-k}T_p^*M\)
由对任意 \(k\)-形式 \(\alpha,\beta\) 成立的等式
\[
\alpha\wedge *\beta
=\langle\alpha,\beta\rangle\,\mathrm{vol}_g
\]
唯一确定。
\end{definition}

存在性并非显然。对基形式 \(e^I\)，令 \(I^c\) 是补指标组，
定义 \(*e^I=\operatorname{sgn}(I,I^c)e^{I^c}\)，
其中符号由
\(e^I\wedge e^{I^c}
=\operatorname{sgn}(I,I^c)\mathrm{vol}_g\) 决定。
这样 \(e^J\wedge *e^I=\delta_{IJ}\mathrm{vol}_g\)，
再由双线性即可得到定义式。若有另一算子满足同式，
它与此算子之差同所有基 \(e^I\) 作楔积都为零；
补基系数逐一为零，故唯一。

在 Euclid 三维的正向直角坐标下，
\[
*1=\dd x\wedge\dd y\wedge\dd z,
\quad *\dd x=\dd y\wedge\dd z,
\quad *(\dd x\wedge\dd y)=\dd z,
\quad *(\dd x\wedge\dd y\wedge\dd z)=1.
\]
再作一次星号时，两组指标交换经过 \(k(n-k)\) 次换位，因而
\begin{equation}
*^2\alpha=(-1)^{k(n-k)}\alpha
\qquad\bigl(\alpha\in\Lambda^kT_p^*M\bigr).
\label{eq:math-hodge-square}
\end{equation}
这条公式以正定度量为前提。若度量签名有 \(q\) 个负方向，
在常用的正向伪正交基约定下还要乘 \((-1)^q\)；
因此只说 \(*^2=\pm1\) 而不说明维数、形式次数和签名是不完整的。

\section{余微分与最熟悉的 Laplace 算子}

对紧支撑的形式，定义全空间内积为
\((\alpha,\beta)=\int_M\alpha\wedge *\beta\)。
在无边界流形，或使边界项消失的边界条件下，
外微分 \(\dd:\Omega^{k-1}(M)\to\Omega^k(M)\)
的形式伴随记为 \(\delta:\Omega^k(M)\to\Omega^{k-1}(M)\)，即
\((\dd\alpha,\beta)=(\alpha,\delta\beta)\)。
先取支集包含在一张坐标图内的形式；普通的逐坐标分部积分
使全微分的积分为零。所得等式只涉及点处运算，
故在重叠坐标图上也相同。
注意乘积求导式
\(\dd(\alpha\wedge *\beta)
=\dd\alpha\wedge *\beta+(-1)^{k-1}\alpha\wedge\dd(*\beta)\)，
把全微分项的积分移到另一边，再用星号平方公式，得到
\[
\delta\beta=(-1)^{n(k+1)+1}*\dd *\beta
\quad\text{在 Riemann 度量下}.
\]
边界不消失时，必须额外写出边界积分，不能把形式伴随
直接称为具体 Hilbert 空间算子的伴随。

Laplace--de Rham 算子定义为
\(\Delta_{\rm dR}=\dd\delta+\delta\dd\)。
对零形式即函数 \(f\)，\(\delta f=0\)，故
\(\Delta_{\rm dR}f=\delta\dd f\)。
在 Euclid 坐标中
\(\dd f=\sum_i(\partial_i f)\dd x^i\)，
由上述伴随公式逐项算得
\[
\Delta_{\rm dR}f=-\sum_i\partial_i^2f.
\]
它是分析中非负的 Laplacian \(-\nabla^2\)，与物理教材
常写的 \(\nabla^2\) 只差一个符号约定。
这一例子同时说明，抽象的形式语言最终必须能退回熟悉的偏微分算子。

极坐标不是另一套定义，只是同一算子的另一种写法。
由正交余标架 \(e^1=dr,e^2=r\,d\varphi\) 有
\(*dr=r\,d\varphi\)、\(*d\varphi=-dr/r\)。因此
\[
*df=r(\partial_r f)\,d\varphi
-\frac{\partial_\varphi f}{r}\,dr,
\quad
d*df=\left[\partial_r(r\partial_r f)
+\frac1r\partial_\varphi^2f\right]dr\wedge d\varphi.
\]
因为体积形式是 \(r\,dr\wedge d\varphi\)，再用
\(\delta df=-*d*df\)，得
\[
\Delta_{\rm dR}f
=-\frac1r\partial_r(r\partial_r f)
-\frac1{r^2}\partial_\varphi^2f.
\]
如果误把 \(d\varphi\) 当成单位长度的一形式，
第二项就会少一个 \(r^{-2}\)；余标架的长度不能省略。

电磁学给出另一个直接检验。在 Euclid 三维，先用度量把向量场
\(V\) 变成一形式 \(V^\flat=g(V,\cdot)\)；
把磁场向量 \(\mathbf B\) 记成二形式
\(\mathcal B=*\mathbf B^\flat\)。
按前面的星号表，
\[
\mathcal B=B_x\,dy\wedge dz
+B_y\,dz\wedge dx+B_z\,dx\wedge dy,
\qquad
d\mathcal B=(\nabla\cdot\mathbf B)\,dx\wedge dy\wedge dz.
\]
对电场一形式 \(\mathcal E=\mathbf E^\flat\) 同样逐项求导，
\(d\mathcal E=*(\nabla\times\mathbf E)^\flat\)。
所以 Maxwell 的两条无源运动学方程
\(\nabla\cdot\mathbf B=0\)、
\(\nabla\times\mathbf E=-\partial_t\mathbf B\)
恰可写成
\[
d\mathcal B=0,
\qquad
d\mathcal E+\partial_t\mathcal B=0.
\]
外微分的 \(d^2=0\) 还解释了为何若局部写
\(\mathcal B=d\mathcal A\)，第一条方程自动成立；
上章的穿孔平面反例同时提醒，局部势的存在不自动保证
它能作为全局单值对象使用。

\begin{exercise}
在 Euclid 二维的正向坐标中求 \(*\dd x\) 与 \(*\dd y\)，
并直接检验式~\eqref{eq:math-hodge-square} 对一形式给 \(*^2=-1\)。
\end{exercise}

\begin{exercise}
对单位球面度量
\(g=d\vartheta^2+\sin^2\vartheta\,d\varphi^2\)，
用同一办法计算 \(\Delta_{\rm dR}f\)。
特别检验 \(f=\cos\vartheta\) 满足
\(\Delta_{\rm dR}f=2f\)。
\end{exercise}

